\chapter{职业生理学、心理学与工效学}
\setcounter{chapter}{2}
\chapterintro{
\textbf{【教学大纲要求】}

\imp{掌握：}体力劳动时的能量代谢、氧消耗的动态；劳动强度分级；动力定型；脑力劳动的生理特点与卫生要求；劳动类型与作业类型（静力作业的特点\keyword{*}）；劳动负荷评价；职业紧张的概念、职业紧张模式（JD-C、ERI）、劳动过程中的紧张源；劳动过程中作业能力的动态变化；影响作业能力的主要因素；人类工效学的概念、意义与内容（人、机器与作业环境）；劳动过程中引起的疾病及其预防（强迫体位、个别器官紧张、压迫及摩擦）。

\imp{熟悉：}职业活动生理检查方法；工作中的心理与社会因素；心身疾病；提高作业能力的主要措施；工作过程的生物力学（肌肉骨骼力学特性、合理用力）；人体测量与应用；机器和工作环境。

（注：\keyword{*} 标示教学难点）
}

% ----- 2.1 职业生理学 -----
\section{职业生理学基础}

\subsection{能量代谢}

\begin{defbox}
\textbf{职业生理学（Work Physiology）}：研究一定劳动条件下的器官和系统功能。
\textbf{职业工效学（Ergonomics）}：研究人-机-环境系统的协调关系。
\textbf{职业心理学（Industrial Psychology）}：研究人的劳动行为与劳动条件之间的关系。
\end{defbox}

\begin{defbox}
\textbf{能量代谢（Energy Metabolism）}：物质代谢过程中的能量释放、转移和利用。

\textbf{基础代谢}：维持基本生命活动所需的最低能量消耗。包括合成代谢和安静代谢（机体状态、睡眠代谢、劳动代谢、分解代谢、食物特殊动力学效应）。

\textbf{氧需（Oxygen Demand）}：劳动1分钟所需的氧量。

\textbf{最大摄氧量（Maximum Oxygen Uptake）}：血液在1分钟内能供应的最大氧量（氧上限）。正常成年人约3 L，运动员约4 L。

\textbf{氧债（Oxygen Debt）}：氧需和实际供氧量之差。
\end{defbox}

\begin{keybox}
\textbf{两种作业类型的氧债特征：}

\begin{table}[H]
\centering
\caption{非乳酸氧债与乳酸氧债的比较}
\begin{tabularx}{\textwidth}{lXX}
\hline
\textbf{特征} & \textbf{非乳酸氧债（较轻劳动）} & \textbf{乳酸氧债（较重劳动）} \\
\hline
作业初时氧来源 & 呼吸、血氧、氧储备 & 呼吸、血氧、氧储备 \\
稳定状态 & 氧债恒定 & 氧债持续增加 \\
稳定状态氧来源 & 呼吸、血氧 & 糖原 \\
偿还氧债时间 & 休息2--3 min & 休息数十分钟至1h以上 \\
氧债与偿还关系 & 氧债 = 偿还氧债 & 氧债 < 偿还氧债 \\
\hline
\end{tabularx}
\end{table}
\end{keybox}

\subsection{劳动强度分级}

\begin{keybox}
\textbf{体力劳动强度分级（三级）：}

\begin{table}[H]
\centering
\caption{劳动强度分级}
\begin{tabularx}{\textwidth}{lXXX}
\hline
\textbf{分级} & \textbf{能消耗状态} & \textbf{氧债} & \textbf{持续时间} \\
\hline
\imp{中等强度} & 氧需 $\leq$ 氧上限 & 氧债恒定 & 长时间（我国多数工农业劳动）\\
\imp{大强度} & 氧需 $>$ 氧上限 & 氧债蓄积 & 数分钟至10 min（重件手工锻打）\\
\imp{极大强度} & 氧需 $\approx$ 氧债 & 氧需 $\approx$ 氧债 & 不超过2 min（短跑、游泳比赛）\\
\hline
\end{tabularx}
\end{table}

\textbf{国际劳动组织（ILO）分级指标：}
\begin{itemize}[leftmargin=*]
    \item 耗氧量 (L/min)、能消耗量 (kJ/min)、心率（次/min）
    \item 直肠温度（℃）、排汗量 (mL/h)
    \item 适用范围：氧需 < 氧上限（稳定状态下）
\end{itemize}

\textbf{我国体力劳动强度分级（GBZ 2.2）：}

劳动强度指数 $I = 3T + 7M$

\begin{itemize}[leftmargin=*]
    \item T：劳动时间率 = 工作日净劳动时间(min) / 工作日总工时(min)
    \item M：8h工作日能量代谢率 [kcal/(min·m$^2$)]
    \item 3和7分别为劳动时间率和能量代谢率的计算系数
\end{itemize}
\end{keybox}

\subsection{体力劳动时的机体变化}

\begin{keybox}
\textbf{一、神经系统——动力定型}

\begin{defbox}
\textbf{动力定型（Dynamic Stereotype）}：长期在同一劳动环境中从事某一作业活动时，通过复合条件反射逐渐形成的作业能力提高。长期脱离该作业活动，会破坏已形成的动力定型。
\end{defbox}

\textbf{二、心血管系统}

\begin{table}[H]
\centering
\caption{体力劳动时心血管系统变化}
\begin{tabularx}{\textwidth}{lXXX}
\hline
\textbf{指标} & \textbf{安静时} & \textbf{无锻炼者} & \textbf{有锻炼者} \\
\hline
心率 HR (min$^{-1}$) & 70--80 & 150--200 & 200--260 \\
每搏输出量 SV (mL) & 40--70 & 增加不多 & 150--200 \\
心输出量 COP (L/min) & 3--5 & 20 & 35 \\
\hline
\end{tabularx}
\end{table}

\begin{itemize}[leftmargin=*]
    \item 血压变化：收缩压↑，脉压差↑；运动后期脉压差↓提示糖原储备↓
    \item 血液再分配：血液从内脏流向肌肉
    \item 血液成分：血糖↓、血乳酸↑
    \begin{itemize}
        \item 安静时血乳酸浓度：1.0 mmol/L
        \item 中等体力劳动：2.3 mmol/L
        \item 重度体力劳动：4.0 mmol/L
        \item 极重体力劳动：15.0 mmol/L
    \end{itemize}
\end{itemize}

\textbf{三、呼吸系统}
\begin{itemize}[leftmargin=*]
    \item 肺通气量增加（呼吸次数↑、肺活量↑）
    \item \keyword{决定最大摄氧量的主要因素是心血管系统的功能}（而非呼吸系统）
\end{itemize}
\end{keybox}

\subsection{脑力劳动的特点}

\begin{itemize}[leftmargin=*]
    \item 特点：抽象性、创造性、非重复性
    \item \imp{智力结构}：一般智力（观察力、记忆力、注意力、思考力、想象力）+ 特殊智力（特殊活动领域内发生作用的能力）
    \item 脑的相对氧代谢高，为等量肌肉的\imp{15--20倍}，\imp{葡萄糖}是最重要的能量来源
    \item 脑力劳动总能量消耗不高，最紧张的脑力劳动能耗不超过基础代谢的10\%
    \item 脑力劳动可使心率减慢，但紧张劳动时心率加快、血压上升、呼吸加速
    \item 血液成分、尿量、尿液成分、汗液、体温等变化不明显
\end{itemize}

\subsection{劳动类型的分类}

\begin{table}[H]
\centering
\caption{劳动类型的分类}
\begin{tabularx}{\textwidth}{lllX}
\hline
\textbf{劳动种类} & \textbf{劳动形式} & \textbf{涉及器官} & \textbf{举例} \\
\hline
\imp{能量性劳动} & 肌力式劳动（付出体力） & 肌肉、骨骼、循环、呼吸 & 搬运、铲砂子 \\
\imp{感觉运动式劳动} & 手和臂精确活动 & 肌肉、肌腱、感官 & 流水线装配、驾驶 \\
\imp{信息性劳动（反应式）} & 吸收加工信息 & 感官（肌肉） & 警卫、监控 \\
\imp{信息性劳动（综合式）} & 信息转换输出 & 感官、脑力 & 编程、翻译 \\
\imp{信息性劳动（创造式）} & 产生信息输出 & 脑力 & 发明、解决问题 \\
\hline
\end{tabularx}
\end{table}

\subsection{静力作业与动力作业}

\begin{keybox}
\textbf{静力作业（Static Work）vs 动力作业（Dynamic Work）：}

\begin{table}[H]
\centering
\caption{静力作业与动力作业的比较}
\begin{tabularx}{\textwidth}{lXX}
\hline
\textbf{对比项} & \textbf{静力（态）作业} & \textbf{动力（态）作业} \\
\hline
肌肉收缩 & 等长收缩 & 等张收缩 \\
关节是否活动 & 否 & 是 \\
能量消耗 & 不高 & 高 \\
血液灌流 & 不充分（压迫血管） & 充分 \\
疲劳程度 & 容易疲劳 & 不易疲劳 \\
\hline
\end{tabularx}
\end{table}
\end{keybox}

% ----- 2.2 劳动负荷与应激 -----
\section{劳动负荷评价}

\begin{defbox}
\textbf{劳动系统（Work System）}：劳动者、劳动对象、劳动工具、劳动环境、产品等相互作用的元素构成的整体。

\textbf{负荷（Stress）}：劳动系统对机体生理心理总的需求和产生的压力，强调外界因素和情景。

\textbf{应激（Strain）}：负荷对具体个人的影响，强调在负荷作用下机体内部的生物过程和反应。

\textbf{适宜水平}：在该负荷下能够连续劳动8小时，不至于疲劳，长期劳动时也不损害健康的卫生限值。约为最大摄氧量的1/3（男：1.1 L/min，女：0.8 L/min）。
\end{defbox}

\begin{tipbox}
\textbf{劳动负荷的评价方法与指标：}

\textbf{客观方法：}
\begin{itemize}[leftmargin=*]
    \item 体力劳动：能量代谢、心率、肌电、皮温、中心体温、血乳酸
    \item 脑力劳动：瞳孔测量、心率与心率变异性、脑诱发电位
\end{itemize}

\textbf{主观方法：}
\begin{itemize}[leftmargin=*]
    \item 体力劳动：调查表、访谈
    \item 脑力劳动：Cooper-Harper量表、SWAT评估法、NASA任务负荷指数
\end{itemize}

\textbf{观察方法：}AET工作分析法、OWAS观察分析劳动姿势和负荷大小
\end{tipbox}

% ----- 2.3 职业心理学 -----
\section{职业心理学与职业紧张}

\subsection{职业心理因素}

\begin{itemize}[leftmargin=*]
    \item \imp{单调作业（Monotonous Work）}
    \item \imp{夜班作业（Night Work）}
    \item 物理因素作业、生产性毒物作业、粉尘作业、脑力作业
\end{itemize}

\subsection{职业紧张模型}

\begin{keybox}
\textbf{职业紧张（Occupational Stress）}：在某种职业条件下，客观需求与个人适应能力之间的失衡所带来的生理和心理压力。

\textbf{紧张反应曲线：}平静（Calm）→ 良性紧张（Eustress，绩效提升）→ 不良紧张（Distress，绩效下降）

\textbf{核心概念：}紧张源（Stressor）→ 个体资源（Personal Resources）→ 应对（Modifier）→ 紧张反应（Strain）

\textbf{JD-C模型（工作要求-自主程度模型，Karasek, 1979）：}
\begin{itemize}[leftmargin=*]
    \item A. 高要求-低控制：\textbf{高紧张工作}，易导致职业紧张和健康问题
    \item B. 高要求-高控制：\textbf{积极工作}，可能带来挑战和成长
    \item C. 低要求-高控制：\textbf{低紧张工作}，较为轻松
    \item D. 低要求-低控制：\textbf{被动工作}，可能导致倦怠
\end{itemize}

\textbf{付出-回报失衡模型（ERI, Siegrist, 1996）：}
\begin{itemize}[leftmargin=*]
    \item A. 高努力-低回报：\textbf{高紧张状态}，易导致职业紧张和健康问题
    \item B. 高努力-高回报：\textbf{平衡状态}，可能带来满足感
    \item C. 低努力-高回报：\textbf{低紧张状态}，较为轻松
    \item D. 低努力-低回报：可能导致\textbf{职业倦怠}
\end{itemize}

\textbf{应对资源：}
\begin{itemize}[leftmargin=*]
    \item 个人应对能力：克服困难、认真对待
    \item 社会支持：有助于缓解紧张，提高职业满意度和职业生命质量
\end{itemize}
\end{keybox}

\subsection{职业紧张反应的表现}

\begin{keybox}
\textbf{心理反应：}工作不满意、躯体不适、疲倦感、抑郁、注意力不集中、易怒等

\textbf{生理反应：}血压升高、心率加快、皮肤生物电反应增强；血与尿中儿茶酚胺和17-羟类固醇增加

\textbf{行为表现：}吸烟、酗酒、频繁就医、药物依赖、不愿参加集体活动等

\textbf{精疲力竭（Burnout）：}
\begin{itemize}[leftmargin=*]
    \item 核心表现：情绪耗竭、人格解体、职业效能下降
    \item 数据：91.1\%中国员工出现疲劳、倦怠、压力等症状，15.4\%出现7种以上症状，5\%出现10种以上症状
\end{itemize}
\end{keybox}

\subsection{心身疾病}

\begin{defbox}
\textbf{心身疾病（Psychophysiological Disorder）}：又称心身障碍或心理生理障碍，是一组与心理和社会因素密切相关，但以躯体症状表现为主的疾病。

常见心身疾病：支气管哮喘、消化性溃疡、原发性高血压、癌症、甲状腺功能亢进。
\end{defbox}

% ----- 2.4 工效学 -----
\section{工效学原理与应用}

\begin{defbox}
\textbf{工效学（Ergonomics）}：研究人-机-环境系统中人的生理、心理特征与机器、环境的协调关系，使劳动工具、劳动环境和劳动组织适合人的生理和心理特征，以达到安全、健康、舒适和高效的目的。
\end{defbox}

\subsection{劳动过程中的工效学问题}

\begin{itemize}[leftmargin=*]
    \item \imp{强制体位引起的疾病}：扁平足、下肢静脉曲张、脊柱弯曲、腰背痛
    \item \imp{个别器官过度紧张}：视力紧张（近视）、发音器官紧张（声带小结）
    \item \imp{压迫及摩擦引起的疾病}：胼胝、滑囊炎
\end{itemize}

\subsection{生物力学应用}

\begin{itemize}[leftmargin=*]
    \item \imp{人体测量学}：工作台、座椅设计应符合人体尺寸
    \item \imp{生物力学}：合理用力避免肌肉骨骼损伤
    \item \imp{人机界面设计}：显示器、控制器的设计应符合人的认知特征
\end{itemize}

% ----- 2.5 影响作业能力的主要因素 -----
\section{影响作业能力的主要因素}

\begin{enumerate}[leftmargin=*]
    \item \imp{社会与心理因素}：工作动机、人际关系、组织氛围
    \item \imp{个体因素}：年龄、性别、体质、健康状况、营养、技能熟练度
    \item \imp{环境因素}：气温、噪声、照明、有害物质等
    \item \imp{劳动组织}：工作时间、休息制度、轮班制度
    \item \imp{工效学因素}：工具设计的合理性、人机匹配程度
\end{enumerate}

\subsection{作业能力的动态变化}

\begin{itemize}[leftmargin=*]
    \item \imp{工作入门期}：作业开始后1--2小时，效率较低
    \item \imp{稳定期}：工作能力稳定，效率最高
    \item \imp{疲劳期}：工作能力下降，出现疲劳感
    \item \imp{终末激发}：临近下班时出现的短暂工作效率提高
\end{itemize}

% ----- 2.6 复习题 -----
\section{复习题}

\begin{enumerate}[leftmargin=*, itemsep=8pt]
    \item \textbf{【名词解释】}能量代谢；氧需；最大摄氧量；氧债；动力定型
    \item \textbf{【名词解释】}静力作业与动力作业；劳动系统；负荷与应激
    \item \textbf{【名词解释】}职业紧张；精疲力竭（Burnout）；心身疾病
    \item \textbf{【名词解释】}工效学（Ergonomics）
    \item \textbf{【简答题】}简述体力劳动强度分级（三级）及各级特征。
    \item \textbf{【简答题】}简述静力作业与动力作业的主要区别。
    \item \textbf{【简答题】}简述JD-C模型和付出-回报失衡模型的主要内容。
    \item \textbf{【简答题】}简述职业紧张反应的表现（心理、生理、行为三个层面）。
    \item \textbf{【简答题】}简述影响作业能力的主要因素。
    \item \textbf{【简答题】}简述作业能力的动态变化过程。
\end{enumerate}

\subsection*{参考答案要点}

\begin{enumerate}[leftmargin=*, itemsep=5pt]
    \item \textbf{能量代谢}：物质代谢过程中的能量释放、转移和利用。\textbf{氧需}：劳动1分钟所需的氧量。\textbf{最大摄氧量}：血液在1分钟内能供应的最大氧量（约3L）。\textbf{氧债}：氧需与实际供氧量之差。\textbf{动力定型}：长期从事同一作业通过复合条件反射形成的作业能力提高。
    \item \textbf{静力作业}：等长收缩，关节不活动，能量消耗不高，血液灌流不充分，容易疲劳。\textbf{动力作业}：等张收缩，关节活动，能量消耗高，血液灌流充分，不易疲劳。\textbf{劳动系统}：劳动者、劳动对象、工具、环境和产品等相互作用的整体。\textbf{负荷}指外界对机体的需求，\textbf{应激}指机体对负荷的反应。
    \item \textbf{职业紧张}：客观需求与个人适应能力之间的失衡所带来的生理和心理压力。\textbf{精疲力竭}：情绪耗竭、人格解体、职业效能下降的综合征。\textbf{心身疾病}：与心理和社会因素密切相关但以躯体症状表现为主的疾病。
    \item \textbf{工效学}：研究人-机-环境系统的协调关系，使劳动工具、环境和劳动组织适合人的生理和心理特征。
    \item 三级：中等强度（氧需≤氧上限，氧债恒定，长时间）、大强度（氧需>氧上限，氧债蓄积，数分钟至10min）、极大强度（氧需≈氧债，不超过2min）。
    \item 五项区别：肌肉收缩方式（等长 vs 等张）、关节活动（否 vs 是）、能量消耗（低 vs 高）、血液灌流（不充分 vs 充分）、疲劳程度（易 vs 不易）。
    \item JD-C模型：按工作要求和自主程度分为四象限（高紧张/积极/低紧张/被动工作）。ERI模型：按努力和回报分为高紧张/平衡/低紧张/倦怠。
    \item 心理：工作不满意、疲倦、抑郁、注意力不集中、易怒。生理：血压↑、心率↑、儿茶酚胺↑。行为：吸烟、酗酒、频繁就医、药物依赖。
    \item 五大因素：社会与心理因素、个体因素（年龄、性别、健康等）、环境因素（气温、噪声、照明等）、劳动组织（工休制度等）、工效学因素（工具设计等）。
    \item 四个阶段：工作入门期（效率较低）→稳定期（效率最高）→疲劳期（效率下降）→终末激发（下班前短暂提高）。
\end{enumerate}

\newpage
% =====================================================================
%  CHAPTER 3: 生产性毒物与职业中毒
% =====================================================================
\newpage